% Logigramme du calcul stationnaire (main.c, newton.c)
% Nécessite dans le préambule : \usetikzlibrary{shapes.geometric}
\begin{figure}[H]
\centering
\begin{tikzpicture}[
    node distance = 0.45cm and 1.0cm,
    >={Stealth[length=2.2mm]},
    etape/.style  = {rectangle, rounded corners=2pt, draw=black!70, fill=blue!6,
                     minimum width=7cm, minimum height=0.75cm, align=center,
                     font=\small},
    debut/.style  = {etape, fill=black!8},
    test/.style   = {diamond, aspect=2.6, draw=black!70, fill=orange!10,
                     inner sep=1pt, align=center, font=\small},
    sortie/.style = {rectangle, rounded corners=2pt, draw=black!70,
                     minimum width=3.2cm, minimum height=0.8cm, align=center,
                     font=\small},
    fleche/.style = {->, thick, draw=black!75},
    etiq/.style   = {font=\footnotesize, fill=white, inner sep=1pt}
]

% ---------------- Préparation ----------------
\node[debut] (lecture) {Lecture des données\\ \texttt{donnees.ini}};
\node[etape, below=of lecture] (maillage) {Construction du maillage\\ nœuds $z_k$, pas $h_k$, volumes $\omega_k$};
\node[etape, below=of maillage] (init) {Initialisation par la solution de Thiele\\ $U^{(0)}$, puis $m = 0$};

% ---------------- Boucle de Newton ----------------
\node[etape, below=0.75cm of init] (assemblage) {Assemblage du résidu $\mathcal{F}(U^{(m)})$\\ et de la jacobienne $J(U^{(m)})$};
\node[etape, below=of assemblage] (echelle) {Mise à l'échelle des lignes de $J$ et de $\mathcal{F}$};
\node[etape, below=of echelle] (resolution) {Résolution de $J\,\delta U = -\mathcal{F}$\\ (Gauss bande, pivot partiel)};
\node[test, below=of resolution] (fini) {Résolution possible\\ et $\delta U$ fini ?};
\node[etape, below=of fini] (maj) {Mise à jour amortie\\ $U^{(m+1)} = U^{(m)} + \delta U$ amorti};
\node[test, below=of maj] (conv) {$\max|\delta C| < \mathrm{tol}_C$\\ et $\max|\delta\varphi| < \mathrm{tol}_\varphi$ ?};
\node[test, below=of conv] (iter) {$m + 1 < m_{\max}$ ?};

% ---------------- Sorties ----------------
\node[sortie, fill=green!10, right=of conv] (succes) {Écriture des profils\\ \texttt{profils.csv}};
\node[sortie, fill=red!10, right=of fini]   (echec)  {Échec\\ pas de convergence};

% ---------------- Flèches principales ----------------
\draw[fleche] (lecture)    -- (maillage);
\draw[fleche] (maillage)   -- (init);
\draw[fleche] (init)       -- (assemblage);
\draw[fleche] (assemblage) -- (echelle);
\draw[fleche] (echelle)    -- (resolution);
\draw[fleche] (resolution) -- (fini);
\draw[fleche] (fini)       -- node[etiq, right] {oui} (maj);
\draw[fleche] (maj)        -- (conv);
\draw[fleche] (conv)       -- node[etiq, right] {non} (iter);

% ---------------- Sorties ----------------
\draw[fleche] (conv) -- node[etiq, above] {oui} (succes);
\draw[fleche] (fini) -- node[etiq, above] {non} (echec);
\coordinate (cote) at ([xshift=0.7cm]succes.east);
\draw[fleche] (iter.east) -- node[etiq, pos=0.5, below] {non} (iter.east -| cote)
              -- (echec.east -| cote) -- (echec.east);

% ---------------- Retour de boucle ----------------
\draw[fleche] (iter.west) -- ++(-1.6,0) node[etiq, pos=0.5, below] {oui}
              |- node[etiq, pos=0.25, left] {$m \leftarrow m+1$} (assemblage.west);

\end{tikzpicture}
\caption{Déroulement du calcul stationnaire. Après la préparation (lecture des
données, maillage, initialisation par la solution de Thiele), la méthode de
Newton répète assemblage, résolution et mise à jour amortie jusqu'à ce que les
corrections sur les concentrations et sur les potentiels passent sous leurs
tolérances respectives.}
\label{fig:logigramme_newton}
\end{figure}
